\shorthandoff{<>}% Babel español trata como abreviaturas los delimitadores de curvatura de Xy-pic.

3) 3. Reducción al caso en que $X$ es liso y afín sobre~$S$.
\marginpar{402}

Según 3)\kern2pt2, puede suponerse que $X$ es liso sobre~$S$. Sea
$(U_i)_{i\in I}$ un recubrimiento finito de~$X$ por abiertos afines, y
sea $X_1$ la suma disjunta de los $U_i$ y
$p\colon X_1 \to X$ el morfismo canónico. Como $p$ es un morfismo de
descenso efectivo para la categoría de los haces étales de tipo finito
sobre esquemas variables, se ve como en 3)\kern2pt2 que, si se supone
demostrado el teorema para los $U_i$, \ie para $X_1$, también vale
para~$X$.

3) 4. Caso en que $X$ es liso y afín sobre~$S$.

Como en 3)1, se ve que, tras restringir $S$ a un abierto no vacío y
efectuar una extensión étale sobreyectiva y una extensión radicial
sobreyectiva de~$S$, puede encontrarse un diagrama conmutativo
$$
\xymatrix{ X \ar[r]^{i} \ar[d]_{f} & P \ar[ld]^{q} \\ Y }
$$
donde $P$ es un esquema liso sobre~$S$, $X$ es el complementario en~$P$
de un divisor de cruzamientos normales relativamente a~$S$
\ifthenelse{\boolean{orig}}
{y $g$ es un morfismo propio.}
{y $q$ es un morfismo propio.}
Si $F$ es un haz constante de orden primo con las características
residuales de~$S$, todo torsor bajo~$F$ está moderadamente ramificado
sobre~$P$ relativamente a~$S$ (\Ref{XIII.2.3}~b)). El hecho de que
$(F,f)$ sea cohomológicamente propio relativamente a~$S$ en dimensión
$\leqslant 1$ y que $f_*\Phi$ sea constructible resulta entonces
de~\Ref{XIII.2.7}.
